\documentclass[12pt, letterpaper]{article}

\title{CS302: Error Handling}
\author{Amiel L. Iliesi}
\date{2026-03-30}

% packages
\usepackage[hidelinks]{hyperref}
\usepackage[x11names]{xcolor}

% new commands
\newcommand{\question}[2][4em]{\textbf{#2}\vspace{#1}\par}
\newcommand{\answer}[1]{\color{Green4}\textbf{#1}\normalcolor}

% new environments
\newenvironment{answerblock}{\color{Green4}}{\normalcolor\vspace{1em}}

% global settings
\setlength{\parindent}{0pt}
\setlength{\parskip}{0.5em}

\begin{document}

\maketitle

\tableofcontents

\pagebreak
\section{C Style: Error Codes}

\question{How do C-style programs and functions typically indicate errors?}

\question{What are some pros and cons to this approach? Why and how does modern
error handling differ?}

\pagebreak
\section{C++}

\subsection{\texttt{stdexcept} and \texttt{throw}}

\question{What is \texttt{stdexcept}? Give a few examples.}

\question{What does the \texttt{throw} keyword do? How is this different from
\texttt{return}?}

\question{What if \texttt{stdexcept} doesn't have what we need? How would you
write something more specific to your needs?}

\subsection{\texttt{try}/\texttt{catch}}

\question{What are \texttt{try} and \texttt{catch}?}

\question{What are the \emph{three} parameter formats you can specify for
\texttt{catch}?}

\subsection{Errors as Return Values, and \texttt{std::expected}}

\textbf{NOTE: you may not need to know this for class, but it's good to know
for modern error handling practices.}

\question{What is an \emph{error as return value}? Describe the pros and cons
to this approach. Also show a simple example of {\lq{throwing}\rq}  an error value with
\texttt{std::expected}.}

\pagebreak
\section{Python}

\subsection{Raise}

\question{What are the differences between Python's \texttt{raise}, and C++'s
\newline\texttt{throw}?}

\question{What types of exceptions can you throw in Python? How can you define
more--if needed?}

\subsection{\texttt{try}/\texttt{except}/\texttt{else}/\texttt{finally}}

\question{Explain these four keywords.}

\end{document}
